\chapter{Introduction}
\label{ch:intro}

\section{Background}
\label{sec:intro-background}

Graph neural networks (GNNs) are now a standard backbone for collaborative filtering. LightGCN~\cite{he2020lightgcn} showed that the feature transformations and nonlinearities used in earlier graph recommenders such as NGCF~\cite{wang2019ngcf} are not necessary for strong ranking performance: linear neighborhood aggregation over the user--item graph, followed by layer-wise averaging, is often enough. Its simplicity makes it a strong default, and later self-supervised methods such as SGL~\cite{wu2021sgl} build on the same backbone.

Knowledge-graph-aware recommendation augments collaborative signals with item-side semantics. KGAT~\cite{wang2019kgat} propagates over a unified user--item--entity graph; KGCL~\cite{yang2022kgcl} and MCCLK~\cite{zou2022mcclk} align KG and collaborative views through contrastive learning; KGRec~\cite{yang2023kgrec} further introduces edge-level rationales. These methods are effective when the KG is dense and reliable. This thesis studies the less favorable, but common, setting in which the KG is sparse and noisy.

\section{An Empirical Anomaly under Sparse KG}
\label{sec:intro-observation}

Our starting point is an empirical anomaly. On the review-derived Amazon Books KG used in this thesis, with only $2.4$ aspect edges per item after filtering, all four representative KG-aware methods fall below pure LightGCN. Table~\ref{tab:intro-baseline} reports the NDCG@20 comparison.

\begin{table}[H]
\centering
\caption{Comparison of KG-aware methods and pure LightGCN on NDCG@20.}
\label{tab:intro-baseline}
\begin{tabular}{lc}
\toprule
Method & NDCG@20 \\
\midrule
MCCLK~\cite{zou2022mcclk}      & 0.1067 \\
KGCL~\cite{yang2022kgcl}       & 0.1073 \\
KGAT~\cite{wang2019kgat}       & 0.1079 \\
KGRec~\cite{yang2023kgrec}     & 0.1095 \\
\midrule
Pure LightGCN~\cite{he2020lightgcn} & \textbf{0.1179} \\
\bottomrule
\end{tabular}
\end{table}
\FloatBarrier

The result does not imply that these baselines are generally weak; they perform well on the dense KG benchmarks used in their original papers. Rather, it exposes a shared failure mode under sparse KG: when KG signal is directly fused into the scoring path, KG noise can enter the same gradient flow that produces the collaborative representation.

\section{Sparse KG as a Practical Setting}
\label{sec:intro-sparse-kg}

Sparse KG is not merely a dataset accident. Review-derived KGs inherit their coverage from what users mention, so many items receive few or uneven aspect tags. Cold-start domains, low-resource languages, and niche product verticals often lack curated sources comparable to Freebase or Wikidata. Privacy-constrained domains may also restrict which relational signals can be exposed. In these settings, KG completion is not a guaranteed remedy because it needs seed signal and can introduce additional noise.

For this reason, robustness when side information is unreliable is as important as peak performance when side information is clean. This thesis takes sparse-KG robustness as the main design objective.

\section{Research Questions}
\label{sec:intro-rq}

The empirical anomaly motivates two questions:

\begin{itemize}
  \item \textbf{RQ1}: Why do existing KG-aware methods underperform pure LightGCN under sparse KG?
  \item \textbf{RQ2}: What architecture can exploit KG signal when useful while preventing KG noise from contaminating collaborative filtering?
\end{itemize}

We answer RQ1 by identifying mandatory KG fusion as the shared structural weakness. We answer RQ2 with RA-GARK, which treats the KG as a gateable side channel rather than as a required part of message passing. The model starts close to pure LightGCN and learns to open the KG channel only when doing so reduces ranking loss.

\section{Thesis Organization}
\label{sec:intro-organization}

Chapter~\ref{ch:relatedwork} reviews graph collaborative filtering, KG-aware recommendation, gating mechanisms, and attention normalization. Chapter~\ref{ch:methodology} presents RA-GARK: the LightGCN local view, KG-SVD aspect-slot initialization, softmax rationale masking, local-biased fusion gate, contrastive regularization, and training procedure. Chapter~\ref{ch:experiments} reports results, ablations, sensitivity analyses, stability checks, a case study, and cost. Chapter~\ref{ch:conclusion} summarizes the findings, limitations, and future work.

\section{Contributions}
\label{sec:intro-contributions}

This thesis makes four contributions:

\begin{enumerate}
  \item It documents a sparse-KG failure case in which KGAT, KGCL, MCCLK, and KGRec all underperform pure LightGCN.
  \item It proposes RA-GARK, a dual-view recommender that integrates KG information through a gateable side channel instead of mandatory fusion.
  \item It shows that RA-GARK outperforms both the strongest KG-aware baseline and pure LightGCN on the same sparse KG.
  \item It reports a methodological observation on attention normalization: in this setting, softmax rationale masking is much more stable than sigmoid masking, suggesting that normalization should match the semantics and scale constraints of the auxiliary channel.
\end{enumerate}
